\section{The stopped leading combined baseline}
\label{sec:terminal-modal-stopped-baseline}

The separate early convolution in
\eqref{eq:terminal-modal-half-velocity-finite-open} discards useful entry
correction.  At leading order the correlated quantity has a stronger stopped
comparison.  This section proves that comparison exactly.  It does not yet
bound the finite-$q$ perturbation.

Retain the leading correction response $\ell_h(r)$ from
\eqref{eq:terminal-modal-finite-l-coefficients} and put
\begin{equation}
 w_h(d)=\ell_h(h-d)+2\ell_h(h-d-1)-3\ell_h(h-d-2).
 \label{eq:terminal-modal-stopped-leading-w}
\end{equation}
For $k\ge1$ and $r\ge0$, define
\begin{equation}
 S_{k,r}=\sum_{h=1}^k w_h(k-1+r),
 \qquad D_{k,r}=S_{k,r}-2t_k(r),
 \label{eq:terminal-modal-stopped-leading-SD}
\end{equation}
where $t_k$ is the half-velocity tail in
\eqref{eq:terminal-modal-half-velocity-tail}.

\begin{lemma}[Deleted-source tail dominates the endpoint tail; proved here]
\label{lem:terminal-modal-stopped-tail-domination}
For every $k\ge1$ and $r\ge0$,
\begin{equation}
 S_{k,r}\ge2t_k(r).
 \label{eq:terminal-modal-stopped-tail-domination}
\end{equation}
More precisely,
\begin{align}
 D_{1,0}&=1,\qquad D_{1,1}=\frac12,
 \label{eq:terminal-modal-stopped-tail-bases}\\
 D_{k,0}&=\frac13\left(1-\left(-\frac12\right)^{k-2}\right)
 &&(k\ge2),
 \label{eq:terminal-modal-stopped-tail-axis}\\
 D_{k+1,r}&=\frac12\{D_{k,r-1}+D_{k,r}\}
 &&(k\ge1,\ 1\le r\le k+1),
 \label{eq:terminal-modal-stopped-tail-Pascal}
\end{align}
with coefficients outside the triangular support interpreted as zero.
\end{lemma}

\begin{proof}
Substitute the binomial-tail formula
\eqref{eq:terminal-modal-finite-l-coefficients} in
\eqref{eq:terminal-modal-stopped-leading-w} and telescope the sum in
\eqref{eq:terminal-modal-stopped-leading-SD}.  At $r=0$, the two remaining
geometric boundary atoms give
\eqref{eq:terminal-modal-stopped-tail-bases}--
\eqref{eq:terminal-modal-stopped-tail-axis}.  For $r\ge1$, Pascal's identity
applied to the remaining binomial atoms gives
\eqref{eq:terminal-modal-stopped-tail-Pascal}; the same identity holds for
$t_k(r)$ and therefore for their difference.  The focused verifier performs
this extraction in exact arithmetic before checking the displayed formulas.

The axis values are nonnegative: the case $k=2$ is zero, and for $k\ge3$
the power in \eqref{eq:terminal-modal-stopped-tail-axis} lies in
$[-1/2,1/4]$.  The two direct bases are positive.  Induction through the
Pascal recurrence proves $D_{k,r}\ge0$ on the whole triangle, which is
\eqref{eq:terminal-modal-stopped-tail-domination}.
\end{proof}

The preceding line calculation is not, by itself, a proof for the folded
path: near the antipode the two reflected source branches coalesce.  We keep
that fold exactly.  Write $c_n^\circ$ for the leading correction in
\eqref{eq:terminal-modal-leading-correction-recurrence}, and let
$c_{m+1}^{\mathrm{prop},\circ}$ be its hypothetical degree-two continuation
from face $m$.  On the even $2m$-cycle put
\begin{equation}
 y_{-1}=\{(c_m^\circ,0)\}^{\mathrm{ev}},\qquad
 y_0=(c_{m+1}^{\mathrm{prop},\circ})^{\mathrm{ev}},\qquad
 y_{k+1}=2Ly_k-Ly_{k-1}.
 \label{eq:terminal-modal-stopped-proper-trajectory}
\end{equation}
Thus the stopped proper half-difference is
\begin{equation}
 Q_{m,k}^{\mathrm{prop}}(j)=y_k(j)-\frac12y_{k-1}(j),
 \qquad k\ge1.
 \label{eq:terminal-modal-stopped-proper-half-difference}
\end{equation}
For $N=m+k+1$, $0\le r\le m$, and $j=m-r$, define the exact folded
deletion coefficients
\begin{equation}
 \widehat S_{m,k,r}
 =\frac{Q_{m,k}^{\mathrm{prop}}(j)-d_N^\circ(j)}{U_0},\qquad
 \widehat D_{m,k,r}=\widehat S_{m,k,r}-2t_k(r),\qquad
 U_0=\frac3{40}.
 \label{eq:terminal-modal-stopped-folded-SD}
\end{equation}

\begin{lemma}[Exact folded stopped tail; proved here]
\label{lem:terminal-modal-stopped-folded-tail}
Let $m\ge64$ and $1\le k<m$.  Then
\begin{equation}
 \widehat S_{m,k,r}\ge2t_k(r),\qquad0\le r\le m.
 \label{eq:terminal-modal-stopped-folded-tail}
\end{equation}
\end{lemma}

\begin{proof}
Direct coefficient extraction from
\eqref{eq:terminal-modal-stopped-proper-trajectory} and
\eqref{eq:terminal-modal-leading-correction-fixed-j} gives
\begin{align}
 \widehat D_{m,1,1}&=\frac12,&
 \widehat D_{m,1,r}&=0 &&(2\le r\le m),
 \label{eq:terminal-modal-stopped-folded-bases}\\
 \widehat D_{m,k+1,r}
 &=\frac12\{\widehat D_{m,k,r-1}+\widehat D_{m,k,r}\}
 &&(1\le k<m-1,\ 1\le r\le m),
 \label{eq:terminal-modal-stopped-folded-Pascal}\\
 \widehat D_{m,k,0}
 &=4-\frac{p_k+2h_{m,k}}3+\frac6{2^m}.
 \label{eq:terminal-modal-stopped-folded-axis}
\end{align}
Here $p_1=6$, $p_k=8+4(-1/2)^k$ for $k\ge2$, and $h_{m,k}$ is the
coefficient in \eqref{eq:terminal-modal-leading-correction-fixed-j}.  The
last positive term in \eqref{eq:terminal-modal-stopped-folded-axis} is the
exact reflected-seed correction; omitting it would give a lower bound, not
an equality.  The focused verifier extracts the cycle coefficients before
checking all three identities in exact arithmetic.

For clarity, the extraction is finite.  At $k=1$ the direct values of
$c_m^\circ,c_{m+1}^\circ,c_{m+2}^\circ$ give
\eqref{eq:terminal-modal-stopped-folded-bases}.  For $r\ge1$, substitute
$j=m-r$ in the fixed-row formula and apply
$h_{j+1,k}=(h_{j,k}+h_{j+1,k-1})/2$ together with Pascal's identity for
$t_k(r)$; this is \eqref{eq:terminal-modal-stopped-folded-Pascal}.  At
$r=0$, the same substitution leaves the reflected seed coefficient
$6/2^m$ and gives \eqref{eq:terminal-modal-stopped-folded-axis}.

It remains to bound the axis.  Put
\begin{equation}
 a_i=[t^i](2-t)^{-m}
 =2^{-m-i}\binom{m+i-1}{i},\qquad
 R_n=\sum_{i=0}^n\left(-\frac12\right)^{n-i}a_i,
 \label{eq:terminal-modal-stopped-h-aR}
\end{equation}
with $a_i=0$ for $i<0$.  Since
\begin{equation}
 \frac{A(t)}{1+t/2}
 =10t^2-28t+24-\frac6{1+t/2},
\end{equation}
coefficient extraction with $n=k+1$ gives
\begin{equation}
 h_{m,k+1}-h_{m,k}
 =10a_{n-2}-28a_{n-1}+24a_n-6R_n.
 \label{eq:terminal-modal-stopped-h-difference}
\end{equation}
For $n\le m-1$, the sequence $a_i$ is increasing through $i=n$.  The
alternating remainder after its first three reversed terms begins negative
and has decreasing magnitudes, so
\begin{equation}
 R_n\le a_n-\frac12a_{n-1}+\frac14a_{n-2}.
\end{equation}
Consequently the right side of
\eqref{eq:terminal-modal-stopped-h-difference} is at least
\begin{equation}
 18a_n-25a_{n-1}+\frac{17}{2}a_{n-2}.
 \label{eq:terminal-modal-stopped-h-positive-lower}
\end{equation}
Set $d=m-n-1\ge0$.  Using
\[
 \frac{a_{n-1}}{a_n}=\frac{2n}{m+n-1},\qquad
 \frac{a_{n-2}}{a_n}
 =\frac{4n(n-1)}{(m+n-1)(m+n-2)},
\]
and multiplying \eqref{eq:terminal-modal-stopped-h-positive-lower} by the
positive denominator, its sign is the sign of
\begin{equation}
 2\{3n^2+11nd+9d^2-10n-9d\}>0.
\end{equation}
Indeed the quadratic part is at least $3(n+d)^2$, the linear loss is at
most $10(n+d)$, and $n+d=m-1\ge63$.  Hence $h_{m,k}$ increases through
$k=m-1$.

At that last index, use
\begin{equation}
 \frac{A(t)}{(1-t)(1+t/2)}
 =18-10t+\frac2{1-t}-\frac2{1+t/2}.
\end{equation}
The negative-binomial symmetry gives
$\sum_{i=0}^{m-1}a_i=1/2$; moreover
$a_{m-2}=a_{m-1}=:a$ and the reversed alternating sum obeys
$R_{m-1}\ge a-a/2=a/2$.  Therefore
\begin{equation}
 h_{m,m-1}\le1+7a,\qquad
 a=2^{-(2m-1)}\binom{2m-2}{m-1}.
\end{equation}
The ratio of consecutive values of $a$ is $(2m-1)/(2m)<1$.  At $m=64$,
\begin{equation}
 2^{127}-28\binom{126}{63}
 =1163019259149843951707093891970921728>0.
\end{equation}
Thus $a<1/28$ and $h_{m,k}<5/4$ for every $1\le k<m$.  Also $p_k\le9$.
The axis identity now yields
\begin{equation}
 \widehat D_{m,k,0}
 >4-\frac{9+5/2}{3}=\frac16.
\end{equation}
The bases in \eqref{eq:terminal-modal-stopped-folded-bases}, followed by
the Pascal propagation \eqref{eq:terminal-modal-stopped-folded-Pascal},
prove $\widehat D_{m,k,r}\ge0$ on every row.  This is
\eqref{eq:terminal-modal-stopped-folded-tail}.
\end{proof}

We now add the literal degree-one endpoint group.  Let $Q_k^\circ(j)$ denote
the leading, undamped version of the correlated linear half-retention
expression
\begin{equation}
 -\mathcal H_kc-\mathcal T_kd
 \label{eq:terminal-modal-stopped-leading-Q}
\end{equation}
at path row $j$ after $k$ fixed full-face steps.

\begin{lemma}[Stopped leading combined baseline; proved here]
\label{lem:terminal-modal-stopped-leading-baseline}
Let $m\ge64$ and $1\le k<m$.  Then, on every full-path row,
\begin{equation}
 Q_k^\circ(j)\ge d_{m+k+1}^\circ(j)\ge\frac{31}{320},
 \qquad0\le j\le m.
 \label{eq:terminal-modal-stopped-leading-baseline}
\end{equation}
Here $d_{m+k+1}^\circ(j)$ is evaluated on the continued growing path; every
displayed row is shared because $j\le m\le m+k-1$.
\end{lemma}

\begin{proof}
Put
$\wp=P_m/q^2$.  The exact frontier formula and
$\chi^m<113/128$ give
\begin{equation}
 \frac{3488}{1921}\le\wp\le2.
 \label{eq:terminal-modal-stopped-frontier-two}
\end{equation}
Indeed $2m\operatorname{artanh}q>1/8$ and
$e^{1/8}>1+1/8+1/128+1/3072>128/113$, proving the bound on $\chi^m$.
For the upper bound on $\wp$, after clearing the positive denominator in
\eqref{eq:terminal-modal-frontier-optimum}, the claim is equivalent to
\[
 6-q-5q^2-10t+t^2(4-q-5q^2)\ge0,
 \qquad t=\chi^m.
\]
The left side decreases in both $q$ and $t$; at
$(q,t)=(1/1024,113/128)$ it is
$4940251803/17179869184>0$.  The lower bound is already
\eqref{eq:terminal-modal-frontier-bounds}.

At leading scale the endpoint position-correction and velocity-remainder
atoms are respectively
\begin{equation}
 c_{\rm end}=\frac1{10}-\frac{3\wp}{8},
 \qquad d_{\rm end}=\frac{\wp}{8}-\frac1{10}.
 \label{eq:terminal-modal-stopped-endpoint-atoms}
\end{equation}
The antipode $m$ is self-reflected on the even $2m$-cycle, so this endpoint
atom is placed only once.  Accordingly, $H_k(r)$ and $t_k(r)$ below are the
coefficients of the full symmetric folded kernels acting on that single
cycle atom; they are not one-orientation coefficients and receive no extra
factor two.  Thus the endpoint contribution at distance $r$ is
\begin{equation}
 A H_k(r)-B t_k(r),\qquad
 A=\frac{3\wp}{8}-\frac1{10}\ge0,\qquad
 B=\frac{\wp}{8}-\frac1{10}\le\frac3{20}=2U_0,
 \label{eq:terminal-modal-stopped-endpoint-response}
\end{equation}
where $H_k(r)\ge0$ is the endpoint coefficient of the positive kernel
$\mathcal H_k$.  By definition
\eqref{eq:terminal-modal-stopped-folded-SD}, the exact stopped proper
comparison is
$Q_{m,k}^{\mathrm{prop}}(j)-d_{m+k+1}^\circ(j)
=U_0\widehat S_{m,k,r}$.  Therefore
\begin{align}
 Q_k^\circ(j)-d_{m+k+1}^\circ(j)
 &=U_0\widehat S_{m,k,r}+AH_k(r)-Bt_k(r)\notag\\
 &\ge U_0\{\widehat S_{m,k,r}-2t_k(r)\}+AH_k(r)
      +(2U_0-B)t_k(r)\notag\\
 &\ge0.
\end{align}
There is no additional reflected endpoint term: the antipode is a single
cycle coordinate, and both the fold in $\widehat S$ and the symmetric
kernel coefficients $H_k,t_k$ already include the two orientations.  This
proves the first inequality in
\eqref{eq:terminal-modal-stopped-leading-baseline}.  The second is
\eqref{eq:terminal-modal-final-leading-margin} applied at
$N=m+k+1\ge64$.
\end{proof}

The gain is deliberately correlated.  The theorem neither bounds
$\mathcal T_kd$ separately nor absorbs the positive-$q$ variation, mass, and
literal endpoint remainders.  Those finite-$q$ terms are the remaining
short-window ledger.
